\documentclass[AER]{AEA}

% --- Packages ---
\usepackage{amsmath,amssymb,mathtools}
\usepackage{newtxtext}
\usepackage{newtxmath}
\usepackage{bm}
\usepackage{booktabs}
\usepackage{array}
\usepackage{placeins}
\usepackage{natbib}
\usepackage{enumitem}
\usepackage{url}
\usepackage[hidelinks]{hyperref}
\usepackage{xcolor}
\usepackage{tikz}

\urlstyle{same}

\newcommand{\qed}{\nolinebreak\hfill\ensuremath{\blacksquare}\par}

\newcommand{\noopsort}[1]{}

\DeclareRobustCommand{\AERlink}[1]{}

% --- 定理环境 (AER标准) ---
\newtheorem{proposition}{PROPOSITION}
\newtheorem{lemma}{LEMMA}
\newtheorem{theorem}{THEOREM}
\newtheorem{corollary}{COROLLARY}
\newtheorem{definition}{DEFINITION}
\newtheorem{remark}{REMARK}
\newtheorem{assumption}{ASSUMPTION}

% --- 核心设置 ---
\issueName{THE AMERICAN ECONOMIC RUBBISH}
\draftSpacing{1.5}
\setcounter{page}{34}
\pubMonth{April}
\pubYear{2026}
\pubVolume{1}
\pubIssue{2}

\begin{document}

\title{Diminishing Marginal Returns to Devotion: A Relational-Contract and Contest Theory of ``Tiangou'' (Simping)$^\dagger$}
\shortTitle{Diminishing Returns to Devotion}

\author{Zhuoyuan Li$^*$}
\date{February 2026}

\begin{abstract}
This paper develops a stylized theory of courtship effort in the attention economy,
focusing on the empirical pattern that additional ``devotion''---messages, favors, and
gifts---delivers sharply diminishing incremental reciprocation. We formalize this
phenomenon using two complementary game-theoretic environments. First, we study a
Bayesian acceptance game in which a receiver (the ``target'') has a private reciprocation
cost and values devotion through a concave reciprocity technology, yielding a concave
success probability for the suitor. Second, we extend to an $N$-suitor attention contest
in which receivers allocate scarce attention across competing suitors. In both settings
we characterize equilibrium effort, show that the marginal effect of additional devotion
on reciprocation is decreasing, and identify conditions under which competition produces
socially excessive devotion despite individually diminishing marginal returns. We then
interpret the results as a theory of ``tiangou'' behavior---a suitor whose effective value of
reciprocation is sufficiently high relative to outside options---and derive comparative
statics, welfare implications, and platform-design predictions (e.g., message caps, gift
taxation, and commitment devices).
\noindent (\textit{JEL} C72, C73, D82, D86, L82, Z13)
\end{abstract}

\maketitle

\let\thefootnote\relax\footnotetext{\noindent\makebox[0em][r]{$^*$}Li: School of Devotion, Tiangou University (email: tiangou@\allowbreak tu.edu.cn). I am grateful to all the ``targets'' who provided me with high-frequency exogenous shocks of non-reciprocation, which were essential for the empirical calibration of this model.}
\let\thefootnote\relax\footnotetext{\noindent\makebox[0em][r]{$^\dagger$}Go to \url{https://webofnothing.top/noi/10.N0/aer.202604.A.0101.0005} to visit the article page for additional materials.}

\def\thefootnote{\arabic{footnote}}
\setcounter{footnote}{0}

\begin{tikzpicture}[remember picture, overlay]
  \node[anchor=north west, align=left, font=\scriptsize] at ([xshift=2cm, yshift=-1.5cm]current page.north west) {
    \textit{American Economic Rubbish 2026, 1(2): 34--44} \\
    \href{https://webofnothing.top/noi/10.N0/aer.202604.A.0101.0005}{\textit{https://webofnothing.top/noi/10.N0/aer.202604.A.0101.0005}}
  };
\end{tikzpicture}

% =========================================================
% Introduction (unnumbered)
% =========================================================

Online interaction has made courtship increasingly resemble an attention market: time and
responsiveness are scarce, while messages, favors, and gifts are abundant and scalable.
A pervasive stylized fact in popular discourse is \textit{diminishing marginal
reciprocation}: the first incremental unit of devotion (a thoughtful message, a small
favor) may meaningfully increase the likelihood of a reply, while later increments produce
much smaller changes in reciprocation. In Chinese internet vernacular this dynamic is
often summarized by the term ``tiangou'' (literally ``licking dog''), referring to a
suitor who repeatedly supplies devotion despite limited or delayed
reciprocation.\footnote{We treat the term descriptively as a label for a strategic
type---a high-valuation suitor in an attention game---and avoid any normative claim about
individuals.}

This paper offers a game-theoretic account of diminishing marginal reciprocation in
courtship, consistent with modern economic theory. Our goal is not to moralize about
relationship behavior, but to provide a clean microfoundation connecting (i) concave
valuation of devotion by receivers, (ii) strategic reciprocation under private information
and attention scarcity, and (iii) equilibrium overinvestment under competition. The theory
also generates testable implications relevant to platforms that intermediate interaction
(dating apps, livestream gifting, social messaging).

Our paper is closely related to \citet{Price2026}, who develops an overlapping framework
analyzing contingent courtship strategies and relational contracts in an ``affection
economy.'' While Pan focuses on appropriable attention and the conditions under which
contingent devotion is self-enforcing, we complement that analysis by formalizing the
diminishing-returns structure of the reciprocity technology and characterizing rent
dissipation in multi-suitor competition. Together, these analyses provide complementary
foundations for a theory of strategic devotion in modern attention markets.

\subsection*{Contributions}

We make three contributions.

\textit{First}, we introduce a minimal Bayesian acceptance game where a receiver has a
private cost of reciprocation and values devotion through a \textit{concave reciprocity
technology}. The suitor chooses effort anticipating acceptance. We characterize
equilibrium effort and show that the marginal benefit of additional devotion is decreasing
whenever the reciprocity technology is concave and costs are convex, generating an
interior optimum even for very high valuations.

\textit{Second}, we extend to a multi-suitor attention contest with a concave ``courtship
score'' mapping effort into effective attractiveness. We derive the symmetric Nash
equilibrium in pure strategies under standard conditions (including log-concavity of the
score), and show that: (i) each suitor faces diminishing marginal returns, yet (ii)
aggregate effort can rise with the number of competitors and dissipate surplus.

\textit{Third}, we analyze welfare and platform design. Because devotion expenditures can
resemble rent-seeking, private incentives may exceed social benefits. We identify
policy-like interventions---message caps, gift taxes, attention quotas, and commitment
tools---that can raise welfare by reducing dissipative competition without eliminating
mutually beneficial matching.

\subsection*{Roadmap}

Section~1 reviews related literature. Section~2 presents the baseline acceptance game and
establishes diminishing marginal returns and comparative statics. Section~3 studies the
$N$-suitor contest and derives equilibrium effort and welfare. Section~4 discusses welfare
and platform design. Section~5 offers extensions (heterogeneous types, signaling,
habituation). Section~6 concludes. Proofs are in the Appendix.

% =========================================================
\section{Related Literature}
% =========================================================

Our analysis draws from several literatures.

\textit{Simping, contingent courtship, and the affection economy.} Most closely related is
\citet{Price2026}, who models courtship as a relational contract problem in which attention
is ``appropriable'' and devotion is contingent on reciprocation. Pan establishes conditions
under which contingent strategies are self-enforcing and characterizes the resulting
equilibrium attention allocations. Our paper complements Pan's analysis by introducing a
concave reciprocity technology that generates diminishing marginal returns, and by
embedding single-receiver dynamics into an $N$-suitor contest to study rent dissipation
and platform design.

\textit{Contests and rent-seeking.} The multi-suitor model is closely related to
Tullock-style contests and all-pay competition in which agents expend resources to win a
prize \citep{Tullock1980, Konrad2009}. Diminishing returns in our setting are captured by
a concave score function, connecting to contest success functions with decreasing
effectiveness of effort.

\textit{Bayesian games and acceptance decisions.} The baseline acceptance game resembles
screening with private reservation values: a receiver accepts if the realized utility
crosses a type-specific threshold. Our concave reciprocity technology plays the role of a
diminishing marginal ``quality'' signal, related to classical signaling and screening
frameworks \citep{Spence1973, Riley2001}.

\textit{Repeated interaction and relational incentives.} While our baseline is static and
our contest model is one-shot, the economic intuition is consistent with relational
contracting: responsiveness can be a costly action used to incentivize costly supply
\citep{BakerGibbonsMurphy2002, Levin2003}. In Section~5 we sketch a repeated-game
microfoundation in which ``breadcrumbing'' can arise as an incentive scheme subject to
limited commitment.

\textit{Attention economy.} The notion that attention is scarce and valuable has long been
emphasized in economics and organization theory \citep{Simon1971}. Our modeling choice to
treat attention as a costly resource complements modern platform settings where attention
allocation is mediated by design choices (limits, ranking, gating).

% =========================================================
\section{Baseline Model: The Acceptance Game}
% =========================================================

\subsection{Players, Timing, and Primitives}

There are two players: a \textit{suitor} $S$ and a \textit{receiver} $R$.

\paragraph{Timing.}
\begin{enumerate}
  \item $S$ chooses a devotion level (effort) $e \in \mathbb{R}_+$.
  \item $R$ observes $e$ and chooses whether to reciprocate: $a \in \{0, 1\}$, where
    $a = 1$ denotes reciprocation (reply/date/commitment) and $a = 0$ denotes
    non-reciprocation.
\end{enumerate}

\paragraph{Receiver private information.} $R$ has a privately known reciprocation cost
$k \in [0, \bar{k}]$ drawn from cdf $F$ with pdf $f$ continuous and strictly positive on
$(0, \bar{k})$.

\paragraph{Payoffs.} Suitor payoff:
\begin{equation}
  U_S(e, a) = aV - c(e),
\end{equation}
where $V > 0$ is $S$'s value of reciprocation and $c : \mathbb{R}_+ \to \mathbb{R}_+$ is
continuously differentiable, strictly increasing, strictly convex, and $c(0) = 0$.

Receiver payoff:
\begin{equation}
  U_R(e, a; k) = b(e) + a \cdot r(e) - ak.
\end{equation}

Here $b(e)$ captures benefits from devotion that accrue regardless of reciprocation
(e.g., gifts, favors, attention received), with $b$ increasing and concave, $b(0) = 0$.
The term $r(e)$ captures the \textit{incremental relational benefit} from reciprocating
with this suitor (e.g., perceived care, compatibility, perceived quality), and is assumed
increasing and concave with $r(0) = 0$. The key economic force is the concavity of
$r(\cdot)$: additional devotion yields diminishing increments in the receiver's
willingness to reciprocate. This concavity assumption parallels the ``appropriable
attention'' framework in \citet{Price2026}, where sustained devotion raises a receiver's
outside option at a decreasing rate.

\paragraph{Equilibrium concept.} We study Perfect Bayesian Equilibrium (PBE). Since the
receiver moves after observing $e$, the equilibrium reduces to a Bayes--Nash
characterization with a cutoff strategy.

\subsection{Receiver Best Response and the Cutoff Structure}

Given $e$, the receiver chooses $a = 1$ if and only if
\begin{equation}
  b(e) + r(e) - k \geq b(e) \quad \iff \quad r(e) \geq k.
\end{equation}
Thus, the receiver reciprocates if and only if her cost is below the relational benefit.

\begin{lemma}[Cutoff reciprocation]\label{lem:cutoff}
  For any devotion level $e$, there exists a cutoff $k^*(e) = r(e)$ such that $R$
  reciprocates if and only if $k \leq k^*(e)$. The suitor's success probability is
  \begin{equation}
    p(e) = \Pr(a = 1 \mid e) = F(r(e)).
  \end{equation}
\end{lemma}

Lemma~\ref{lem:cutoff} implies that the suitor's expected payoff is
\begin{equation}
  \Pi(e) \equiv \mathbb{E}[U_S(e, a)] = VF(r(e)) - c(e).
\end{equation}

\subsection{Equilibrium Devotion and Diminishing Marginal Returns}

We now establish existence, uniqueness, and the diminishing marginal returns property.

\begin{assumption}[Regularity]\label{ass:regularity}
  (i) $c$ is $C^2$, strictly convex, and $\lim_{e \to \infty} c'(e) = \infty$.\\
  (ii) $r$ is $C^2$, strictly increasing, concave, and $\lim_{e \to \infty} r(e) = \bar{r} < \bar{k}$.\\
  (iii) $F$ is $C^1$ and strictly increasing on $[0, \bar{k}]$.
\end{assumption}

Assumption~\ref{ass:regularity}(ii) allows for a natural saturation: even extreme
devotion cannot guarantee reciprocation if the receiver's cost is high enough.

\begin{proposition}[Existence and characterization]\label{prop:existence}
  Under Assumption~\ref{ass:regularity}, the suitor's problem~(5) has a (possibly
  boundary) maximizer $e^* \in [0, \infty)$. If an interior solution exists, it satisfies
  the first-order condition
  \begin{equation}
    Vf(r(e^*))r'(e^*) = c'(e^*),
  \end{equation}
  and the second-order condition holds whenever $r''(e) \leq 0$ and $c''(e) > 0$.
\end{proposition}

\paragraph{Diminishing marginal returns.} Define the \textit{marginal reciprocation
value} (MRV) of devotion:
\begin{equation}
  \mathrm{MRV}(e) \equiv \frac{d}{de}\bigl[VF(r(e))\bigr] = Vf(r(e))r'(e).
\end{equation}

\begin{proposition}[Diminishing marginal returns to devotion]\label{prop:dmr}
  Suppose $f$ is weakly decreasing on the relevant range and $r$ is concave. Then
  $\mathrm{MRV}(e)$ is weakly decreasing in $e$. In particular, each additional unit of
  devotion raises the success probability by less than the previous one:
  \begin{equation}
    \frac{d^2}{de^2} F(r(e)) \leq 0.
  \end{equation}
\end{proposition}

Proposition~\ref{prop:dmr} formalizes the idea that ``the first message matters most'':
concavity of $r$ (and mild regularity of $F$) implies concavity of the success
probability. The suitor may still choose high devotion if $V$ is large, but the
incremental impact declines with effort.

\subsection{Comparative Statics and the ``Tiangou'' Type}

We interpret tiangou as a suitor with sufficiently high effective value of reciprocation
relative to costs and outside options. To make this formal, allow the suitor to have an
outside option $\bar{u} \geq 0$ (e.g., alternative matches, leisure). Then the objective
becomes
\begin{equation}
  \Pi(e) = VF(r(e)) - c(e) - \bar{u},
\end{equation}
and the suitor participates only if $\max_{e \geq 0} \Pi(e) \geq 0$.

\begin{proposition}[Value and outside-option effects]\label{prop:compstat}
  Assume an interior optimum and conditions for the implicit function theorem. Then:
  \begin{enumerate}
    \item $e^*$ is increasing in $V$.
    \item $e^*$ is decreasing in any parameter that shifts $c$ upward (higher marginal
      cost).
    \item A higher outside option $\bar{u}$ weakly reduces devotion through the
      participation constraint; when participation binds, the equilibrium may jump from
      positive $e^*$ to $0$.
  \end{enumerate}
\end{proposition}

\begin{remark}[A minimal definition of ``tiangou'']\label{rem:tiangou}
  In this baseline model, a tiangou type can be defined as a parameter region in which
  $V$ is high (or $\bar{u}$ low) such that the chosen devotion $e^*$ is large even though
  $\mathrm{MRV}(e)$ is already small. Formally, fix a threshold $\varepsilon > 0$ and
  define $e_\varepsilon$ by $\mathrm{MRV}(e_\varepsilon) = \varepsilon$. A suitor is
  tiangou-like if $e^* \geq e_\varepsilon$, i.e., he invests in the region of low
  marginal reciprocation value. This parallels \citet{Price2026}'s notion of ``excess
  simping''---continued devotion beyond the relational optimum.
\end{remark}

% =========================================================
\section{Competition: An $N$-Suitor Attention Contest}
% =========================================================

We now introduce competition among multiple suitors. This captures environments with
thick markets (social media, dating apps) where receivers face many potential suitors and
attention is allocated competitively.

\subsection{Environment}

There is one receiver $R$ and $N \geq 2$ suitors indexed by $i \in \{1, \ldots, N\}$.
Suitors simultaneously choose devotion $e_i \geq 0$ at cost $c(e_i)$.

The receiver allocates a single unit of reciprocation (e.g., chooses one suitor to
date/respond to). We model the receiver's choice probabilistically via a contest success
function:
\begin{equation}
  \Pr(i \text{ wins}) = \frac{s(e_i)}{\sum_{j=1}^N s(e_j)},
\end{equation}
where $s : \mathbb{R}_+ \to \mathbb{R}_+$ is the \textit{courtship score}, increasing
with $s(0) = 0$. We assume $s$ is concave to capture diminishing marginal effectiveness
of additional devotion. The ``prize'' for winning is reciprocation value $V$.

Suitor $i$'s expected payoff is
\begin{equation}
  \pi_i(e_i, e_{-i}) = V \cdot \frac{s(e_i)}{\sum_{j=1}^N s(e_j)} - c(e_i).
\end{equation}

\subsection{Symmetric Equilibrium}

We focus on symmetric pure-strategy Nash equilibrium.

\begin{assumption}[Contest regularity]\label{ass:contest}
  $s$ is $C^2$, strictly increasing, concave, and log-concave (i.e., $\log s(\cdot)$ is
  concave). $c$ is as in Assumption~\ref{ass:regularity}(i).
\end{assumption}

\begin{proposition}[Symmetric equilibrium in the concave-score contest]\label{prop:contest}
  Under Assumption~\ref{ass:contest}, there exists a symmetric pure-strategy Nash
  equilibrium $e^c$ satisfying
  \begin{equation}
    c'(e^c) = V \cdot \frac{N-1}{N^2} \cdot \frac{s'(e^c)}{s(e^c)}.
  \end{equation}
  Moreover, if $c'$ is strictly increasing and $\frac{s'}{s}$ is strictly decreasing
  (implied by log-concavity), the symmetric equilibrium is unique.
\end{proposition}

Equation~(12) makes the economics transparent: marginal cost equals marginal benefit,
where marginal benefit is proportional to the \textit{semi-elasticity} $s'/s$. Concavity
and log-concavity imply $s'/s$ declines with effort, producing diminishing marginal
returns.

\begin{corollary}[Competition and the ``arms race'']\label{cor:arms}
  In the symmetric equilibrium, individual effort $e^c$ is increasing in $V$ and (under
  mild additional conditions) can increase with $N$ for relevant ranges. Even when $s$ is
  concave (diminishing effectiveness), adding competitors can raise aggregate devotion
  $Ne^c$ and dissipate surplus.
\end{corollary}

\subsection{A Canonical Functional-Form Example}

To connect with standard contest theory and yield closed forms, suppose
\begin{equation}
  s(e) = e^\gamma, \quad \gamma \in (0, 1),
\end{equation}
capturing diminishing effectiveness ($\gamma < 1$), and
\begin{equation}
  c(e) = \frac{\alpha}{2} e^2, \quad \alpha > 0.
\end{equation}
Then $s'/s = \gamma/e$, and (12) gives
\begin{equation}
  \alpha e^c = V \cdot \frac{N-1}{N^2} \cdot \frac{\gamma}{e^c}
  \quad\implies\quad
  e^c = \sqrt{V \cdot \frac{N-1}{N^2} \cdot \frac{\gamma}{\alpha}}.
\end{equation}
The equilibrium effort increases in $V$ and $\gamma$, decreases in $\alpha$, and exhibits
the usual contest effect in $N$.

\begin{remark}[Diminishing marginal returns in the contest]\label{rem:contest_dmr}
  Even though each suitor's marginal win probability from additional devotion is decreasing
  in $e_i$ (due to $\gamma < 1$ and the denominator effect), equilibrium effort can still
  be substantial when the prize $V$ is large or the market is competitive.
\end{remark}

% =========================================================
\section{Welfare and Platform Design}
% =========================================================

Diminishing marginal returns at the individual level does not guarantee efficiency. In
contests, private incentives ignore the negative externality imposed on rivals through the
denominator $\sum_j s(e_j)$.

\subsection{A Simple Welfare Benchmark}

Define social welfare as the sum of suitor utilities plus the receiver's realized matching
utility net of devotion costs. To avoid imposing a particular ``true'' benefit of
devotion, we consider two polar benchmarks.

\paragraph{Benchmark A (pure rent-seeking).} Devotion is purely wasteful from a social
perspective (e.g., excessive messaging time, monetary gifts with no true surplus), so
welfare is
\begin{equation}
  W_A = V - \sum_{i=1}^N c(e_i),
\end{equation}
since exactly one suitor obtains value $V$ in expectation. In the symmetric equilibrium,
expected welfare is
\begin{equation}
  \mathbb{E}[W_A] = V - Nc(e^c),
\end{equation}
which can decrease in $N$ and $V$ for ranges where $Nc(e^c)$ grows rapidly.

\paragraph{Benchmark B (productive devotion).} Devotion creates real surplus for the
receiver, captured by $B\!\left(\sum_i e_i\right)$ with $B$ concave, so
\begin{equation}
  W_B = V + B\!\left(\sum_{i=1}^N e_i\right) - \sum_{i=1}^N c(e_i).
\end{equation}
Concavity of $B$ implies diminishing social marginal returns, so competition can again
overshoot the welfare-optimal total devotion.

\subsection{Design Implications}

We interpret platform features as instruments that change $c(\cdot)$, cap $e$, or alter
the contest success function.

\begin{proposition}[Message caps and gift taxes]\label{prop:policy}
  Consider the contest model with welfare $W_A$ (pure rent-seeking). Any policy that
  weakly reduces equilibrium effort (e.g., a cap $e_i \leq \bar{e}$ below $e^c$, or a
  per-unit tax that increases marginal cost) weakly increases welfare. Under $W_B$
  (productive but concave benefits), such policies increase welfare whenever equilibrium
  total devotion exceeds the social optimum characterized by $B'(\cdot) = c'(\cdot)$ at
  symmetry.
\end{proposition}

\begin{remark}[Interpretation for tiangou dynamics]\label{rem:platform}
  In many platform settings, the suitor's private $V$ may be amplified by design (streaks,
  read receipts, intermittent rewards), effectively increasing $V$ and shifting equilibrium
  toward high devotion in low-marginal-return regions. Conversely, strengthening outside
  options (improved discovery, fair ranking) can reduce tiangou-like overinvestment. These
  platform-design prescriptions complement the commitment-device analysis in
  \citet{Price2026}.
\end{remark}

% =========================================================
\section{Extensions}
% =========================================================

This section sketches several extensions that preserve the core diminishing-marginal-returns result.

\subsection{Heterogeneous Suitors and Signaling}

Let suitor type $\theta \in \{H, L\}$ affect cost: $c(e; \theta)$ with $c_e(e; H) <
c_e(e; L)$. If the receiver values type directly (e.g., compatibility), then devotion can
serve as a costly signal \citep{Spence1973}. Concavity of $r(e)$ implies that extremely
high devotion becomes progressively less informative at the margin, which can compress
separation and lead to pooling or partial pooling equilibria. A testable implication is
that beyond a certain point, additional gifts or messages do little to update beliefs.

\subsection{Habituation and Attention Inflation}

Introduce a state $h$ capturing habituation (reference dependence): $r(e; h)$ is
increasing and concave in $e$ but decreasing in $h$, and $h$ evolves as $h' = \rho h +
(1 - \rho)e$. Then the marginal effectiveness $r_e(e; h)$ declines both in $e$ and in
past devotion. This produces ``attention inflation'': sustaining the same acceptance
probability requires increasing devotion over time, consistent with observed escalation
in some relationships.

\subsection{Relational Contracting and ``Breadcrumbing''}

A repeated-game version can endogenize intermittent reciprocation. Let the receiver choose
attention intensity $a_t \in [0, 1]$ at convex cost $\phi(a_t)$, while receiving concave
benefit $b(e_t)$. A self-enforcing arrangement can implement $a_t$ as an incentive
instrument for $e_t$. Under standard relational-contract constraints
\citep{BakerGibbonsMurphy2002, Levin2003}, promised attention is bounded by the
receiver's temptation to deviate, and this bound inherits concavity from $b(\cdot)$,
again generating diminishing marginal returns: higher devotion can be rewarded, but only
at a decreasing rate. This mechanism is closely related to the contingent-courtship
equilibria characterized in \citet{Price2026}.

\subsection{Empirical Implications}

Although the paper is theoretical, it yields qualitative predictions.
\begin{enumerate}
  \item \textbf{Concavity of response rates.} Holding match fixed, the probability of
    receiving a reply (or the intensity of reciprocation) should be concave in
    messages/gifts when receiver valuation is concave.
  \item \textbf{Competition amplifies waste.} In thicker markets (larger $N$), individual
    devotion may rise and total devotion rises more, even as marginal returns decline.
  \item \textbf{Outside options reduce overinvestment.} Improvements that raise $\bar{u}$
    (better matching, better visibility) reduce devotion and reduce tiangou-like behavior.
  \item \textbf{Platform features shift $V$.} Intermittent reinforcement (streaks, read
    receipts, partial replies) increases the perceived prize $V$, increasing equilibrium
    devotion and pushing behavior into low-marginal-return regions.
\end{enumerate}

% =========================================================
\section{Conclusion}
% =========================================================

This paper provides a game-theoretic framework for understanding diminishing marginal
reciprocation in courtship devotion and its relationship to ``tiangou'' behavior. In a
Bayesian acceptance game, concavity of a receiver's reciprocity technology yields a
concave success probability and declining marginal returns to devotion. In a competitive
attention contest, each suitor still faces diminishing returns, yet competition can
produce high aggregate devotion and welfare losses analogous to rent dissipation. These
results suggest a coherent economic interpretation of widely discussed social phenomena
and highlight the role of platform design in shaping incentives in attention markets. Our
analysis complements the contingent-courtship and relational-contract approach of
\citet{Price2026}, and taken together these frameworks offer a more complete theory of
strategic devotion in the modern affection economy.

% =========================================================
% REFERENCES
% =========================================================
\DeclareRobustCommand{\AERlink}[1]{\makebox[0pt][r]{\href{#1}{\color{blue}$\blacktriangleright$}\hspace{0.3em}}}

\bibliographystyle{aea}
\bibliography{AERub_tiangou2}

% =========================================================
\appendix
\section{Proofs}
% =========================================================

\subsection{Proof of Proposition~\ref{prop:existence}}

Existence follows from continuity of $\Pi(e)$, $\Pi(0) = VF(0) \geq 0$, and the fact
that $c(e) \to \infty$ as $e \to \infty$ while $VF(r(e)) \leq VF(\bar{r}) < \infty$, so
$\Pi(e) \to -\infty$ as $e \to \infty$. Hence $\Pi$ attains a maximum on $[0, \infty)$.
For an interior maximizer, the FOC~(6) holds. Under concavity of $r$ and strict convexity
of $c$, the SOC is satisfied under standard curvature conditions. \qed

\subsection{Proof of Proposition~\ref{prop:dmr}}

Differentiate $F(r(e))$:
\[
  \frac{d}{de} F(r(e)) = f(r(e))\, r'(e).
\]
Differentiate again:
\[
  \frac{d^2}{de^2} F(r(e)) = f'(r(e))\,(r'(e))^2 + f(r(e))\, r''(e).
\]
If $f'(\cdot) \leq 0$ on the relevant range, $r'(e) \geq 0$, and $r''(e) \leq 0$
(concavity), then both terms are weakly negative, so the second derivative is $\leq 0$.
Multiplying by $V > 0$ preserves the inequality, so MRV is weakly decreasing. \qed

\subsection{Proof of Proposition~\ref{prop:contest}}

Fix $e_{-i}$ and consider $\pi_i(e_i, e_{-i})$ in~(11). For interior $e_i$, the FOC is
\[
  V \cdot \frac{s'(e_i)\,\sum_{j \neq i} s(e_j)}{\left(\sum_{j=1}^N s(e_j)\right)^2}
  = c'(e_i).
\]
In a symmetric equilibrium with $e_j = e$ for all $j \neq i$ and $e_i = e$, we have
$\sum_{j \neq i} s(e) = (N-1)s(e)$ and $\sum_{j=1}^N s(e) = Ns(e)$, yielding
\[
  c'(e) = V \cdot \frac{s'(e)(N-1)s(e)}{(Ns(e))^2}
         = V \cdot \frac{N-1}{N^2} \cdot \frac{s'(e)}{s(e)}.
\]
Existence follows from standard arguments using continuity and boundary behavior;
uniqueness follows if $c'$ is increasing and $s'/s$ is decreasing (log-concavity). \qed

\end{document}
